% chapters/03/evaluasi.tex
\bagian{Asesmen dan Instrumen Penilaian}
\label{sec:asesmen-lingkungan}

\begin{asesmenbox}
\textbf{Kedudukan asesmen bab ini.} Bab ini dinilai melalui Tugas Praktikum 01
yang termasuk komponen \istilah{tugas praktikum dan portofolio} pada RPS,
yaitu komponen berbobot 25\% dari nilai akhir yang mengukur
\CPMK{1} sampai \CPMK{4}. Penilaian dilakukan atas artefak yang tersimpan pada
repositori Git Anda, bukan atas jawaban lisan atau hafalan.

\vspace{0.4em}
\textbf{Yang dikumpulkan.} Repositori Git memuat struktur folder proyek,
skrip pemeriksaan lingkungan, program jendela yang telah disunting, berkas
\code{.gitignore}, berkas \code{README.md}, dan catatan praktikum.
\end{asesmenbox}

\subsection*{Instrumen 1: Rubrik Praktikum 01}

\begin{table}[htbp]
\centering
\caption{Rubrik penilaian Praktikum 01}
\label{tab:rubrik-praktikum-01}
{\setlength{\tabcolsep}{4pt}\footnotesize%
\begin{tabularx}{\textwidth}{@{}L{3.8cm} C{1.5cm} Y@{}}
\toprule
\headrow \textbf{Kriteria} & \textbf{Bobot} & \textbf{Deskripsi capaian} \\
\midrule
Kesesuaian fungsional & 40\% &
Program jendela berjalan tanpa kesalahan, menampilkan seluruh baris informasi
versi, dapat diubah ukurannya melalui baris perintah, dan menutup dengan rapi
ketika jendela ditutup. \\
\addlinespace
Kualitas kode dan struktur & 30\% &
Berkas tersusun pada folder yang benar, penamaan konsisten, terdapat komentar
pada bagian yang perlu dijelaskan, dan tidak ada kode sisa yang tidak dipakai.
Berkas \code{.gitignore} tersedia dan berfungsi. \\
\addlinespace
Kelengkapan dokumentasi & 30\% &
Berkas \code{README.md} menerangkan isi repositori, catatan praktikum memuat
kendala beserta cara penyelesaiannya, dan pesan penyimpanan Git menerangkan
pekerjaan yang disimpan. \\
\bottomrule
\end{tabularx}}
\end{table}

\subsection*{Instrumen 2: Soal Konsep}

Jawablah dengan singkat, tanpa membuka bab terlebih dahulu. Kunci jawaban
disajikan pada Tabel~\ref{tab:kunci-konsep-ch03}.

\begin{enumerate}
  \item Apa pekerjaan yang dilakukan \code{pygame.init()}?
  \item Apa bedanya menggambar dengan \code{screen.blit} dan menayangkan
        dengan \code{pygame.display.flip()}?
  \item Mengapa program pada bab ini belum disebut memiliki siklus permainan?
  \item Mengapa pemasangan Pygame sebaiknya memakai bentuk
        \code{python -m pip}?
\end{enumerate}

\begin{table}[htbp]
\centering
\caption{Kunci soal konsep Bab III}
\label{tab:kunci-konsep-ch03}
{\setlength{\tabcolsep}{4pt}\footnotesize%
\begin{tabularx}{\textwidth}{@{}C{0.8cm} Y@{}}
\toprule
\headrow \textbf{No} & \textbf{Pokok jawaban} \\
\midrule
1 & Menyalakan seluruh bagian pustaka Pygame, yaitu tampilan, huruf, bunyi, dan
penanganan peristiwa, sehingga dapat dipakai oleh program. \\
\addlinespace
2 & Menggambar hanya menyusun isi permukaan di memori, sedangkan menayangkan
memindahkan hasilnya ke layar sehingga terlihat pemain. Tanpa penayangan,
jendela tetap menampilkan gambar lama. \\
\addlinespace
3 & Karena perulangannya hanya menunggu peristiwa. Keadaan permainan tidak
diperbarui dan gambar tidak disusun ulang, sehingga tidak ada perubahan yang
berkelanjutan. \\
\addlinespace
4 & Agar pemasangan mengikuti penafsir Python yang sama dengan yang akan
dipakai menjalankan program, terutama pada komputer yang memiliki lebih dari
satu versi Python. \\
\bottomrule
\end{tabularx}}
\end{table}

\begin{catatabox}
Apabila kriteria kelengkapan dokumentasi belum terpenuhi, lengkapi catatan
praktikum Anda lebih dahulu. Kebiasaan mencatat kendala pada minggu-minggu
awal akan sangat membantu pada Bab~XVI, ketika Anda harus menyusun laporan
pengembangan dan \istilah{log bug} untuk proyek akhir.
\end{catatabox}
